\begin{table}[htbp]
  \centering
  \caption{原协议对象到两槽抽象的映射}
  \label{tab:source-to-abstraction}
  \footnotesize
  \setlength{\tabcolsep}{3pt}
  \begin{tabularx}{\linewidth}{@{}L{0.19\linewidth}L{0.17\linewidth}L{0.23\linewidth}Y@{}}
    \toprule
    原协议对象 & 形式化表示 & 保留语义 & 省略内容与假设 \\
    \midrule
    输入中的公钥、预密钥束及其参与方关联 & $A_i$ & 比较条目所关联的参与方 & 假设参与方投影正确；不验证密钥编码、预密钥及身份解析 \\
    一次发送及其协议消息 & $(\mathit{oid}_i,m_i)$ & 区分发送实例及其新鲜符号消息 & 不表示原协议transcript会话标识，不展开密文结构或密钥派生 \\
    条目的发送来源 & \texttt{!Sent(A,oid,m)} & 处理时要求完整元组精确匹配；来源可被重复读取 & 理想化来源事实，不证明签名、解封装或完整密码学认证 \\
    输入向量$S$中的两个位置 & $(E_1,E_2)$及线性槽位 & 有序收集、攻击者组合和顺序处理 & 机器范围固定为两槽，不推广为任意批长验证 \\
    同次调用与共享接收方临时状态 & $(\mathit{bid},\mathit{rst})$ & 限定共同批次和接收方上下文 & 不是协议线路字段或完整临时密码状态 \\
    按输入分量组织的接收结果 & \texttt{ReceiverAccept}观察事件 & 观察各槽位来源匹配后的抽象接受 & 不产生真实会话标识或密钥，不表示密钥安装；分量输出仅构成动机 \\
    \bottomrule
  \end{tabularx}
  \par\smallskip
  \begin{minipage}{\linewidth}
    \scriptsize 注：表1解释模型内部比较坐标，本表解释原接口到抽象模型的取舍；该映射由研究者建立，尚无协议精化证明。
  \end{minipage}
\end{table}
